\section{The procedure as an admission criterion: several published systems}

Section~6 carried one system through all three steps. The procedure is not a
single reconstruction, however, but a decision rule, and its usefulness is
better shown by what it decides across systems that were measured for other
reasons. Table~\ref{tab:admission} records how far each of several published
metal-carbonyl systems gets, and what stops it. Every input is a value stated in
the running text of its source; none is read off a figure. Values a source gives
only as similar to another number are entered as illustrative, and the source
gate of Section~2.3 then blocks any result derived from them from being quoted
as quantitative, without the authors having to remember to do so.

\begin{table}[htbp]
\centering\small
\noindent\begin{tabular}{@{}p{0.19\linewidth}p{0.09\linewidth}p{0.17\linewidth}p{0.08\linewidth}p{0.09\linewidth}p{0.22\linewidth}@{}}
\toprule
System & \(T\) & Verdict & \(R\) & Over & Limiting evidence \\
\midrule
W(CO)$_6$ / 2-MP & 250--300~K & Step 1 passed & 1 & n/a & no text \(T_1\) at this \(T\) \\
W(CO)$_6$ / DBP & 300~K & Step 1 passed & 26 & n/a & no text \(T_1\) at this \(T\) \\
W(CO)$_6$ / CCl$_4$ & 295~K & Step 1 refused & n/a & 1300 & no echo \\
W(CO)$_6$ / CHCl$_3$ & 295~K & Step 1 refused & n/a & 1200 & no echo \\
Cr(CO)$_6$ / CCl$_4$ & 295~K & Step 1 refused & n/a & 700\(^{*}\) & width inherited as similar \\
Cr(CO)$_6$ / CHCl$_3$ & 295~K & Step 1 refused & n/a & 930\(^{*}\) & width inherited as similar \\
Rh(CO)$_2$acac / DBP & 3.4, 250~K & Steps 1--3 & 136 (3.4~K) & n/a & reference entry (Section~6) \\
\bottomrule
\end{tabular}
\caption{The three-step procedure applied as an admission criterion. Verdict records how far each published system gets and what stops it. \(R\) is the ratio of absorption to echo-isolated homogeneous width; Over is the factor by which the absorption width exceeds the lifetime-limited width, a bound on the total overestimate rather than a decomposition. \(^{*}\)Not reportable: an illustrative input, blocked by the source gate.}
\label{tab:admission}
\end{table}

The two entries at the top are the substantive result, and they concern one
molecule and one mode. The asymmetric CO stretch of W(CO)$_6$ near 1980~cm$^{-1}$
in 2-methylpentane above 250~K has an absorption line that is homogeneously
broadened: the echo signal becomes a free induction decay whose width matches
the measured absorption width \citep{tokmakoff1995}, so \(R_{\mathrm{abs/hom}} = 1\)
and reading that width as a dynamic quantity is correct. The same mode of the
same molecule at 300~K in dibutylphthalate has a homogeneous width near
1~cm$^{-1}$ against an absorption width of 26~cm$^{-1}$ \citep{tokmakoff1995}, so
the absorption width overstates the dynamic width by a factor of about 26 and
must not be used. Solvent alone moves the verdict. Whether a measured linewidth
may be read as dynamics is therefore a property of the system rather than of the
chromophore, and no inspection of the spectrum alone distinguishes the two
cases: the echo measurement is what separates them.

The lower entries record the more common situation, in which no echo exists.
For those, Step~1 cannot be satisfied and the procedure returns a bound rather
than a value. The bound is not weak. Comparing the stated absorption width with
the width implied by the stated lifetime, \(\Gamma_{\mathrm{life}} = 1/2\pi c T_1\),
gives overshoot factors near \(10^3\) for the room-temperature carbonyl solutions
\citep{tokmakoff1994}. Reading such a linewidth as a mechanical damping
coefficient would overestimate \(\chi\) by about three orders of magnitude. That
overshoot lumps inhomogeneous broadening together with pure dephasing and
orientational relaxation, and is therefore a bound on the total overestimate,
not a decomposition of it; separating the three requires the echo the procedure
asks for.

\subsection*{A fourth verdict: when the decomposition itself is not well posed}

The entries above are all metal carbonyls, and all of their refusals have the
same character: a measurement is missing. Applying the procedure outside that
family produces a refusal of a different kind, and the distinction matters
because it is a statement about the system rather than about the data.

For the OH stretch of liquid water the vibrational lifetime is not a single
number. Frequency-resolved relaxation measurements give about 250~fs at the red
edge of the band rising to about 550~fs at the blue edge, with a comparable
factor-of-two variation reported for HOD in D$_2$O \citep{water2015}. Step~2 requires one lifetime in order to define a single
scalar mechanical mode. The local physics is not in question: each hydrogen-bond
environment relaxes perfectly well. What fails is the reduction. A macroscopic
band that represents a dispersed ensemble of local environments is not one
oscillator, so mapping the whole band onto a single scalar \((M, C, K)\) is
ill posed, and there is no single macroscopic mechanical coordinate on which to
place it. 

The OD stretch of HOD in H$_2$O separates the two failures cleanly. Its excited
state decay is wavelength independent with a decay time of 1.45~ps \citep{steinel2004}, so Step~2
would be well posed. Step~1 is not: the source reports that the dynamic line width, its
measure of spectral diffusion, reaches 88\% of the total absorption line width by 0.5~ps,
95\% by 1.0~ps, and 98\% by 1.5~ps, so a hydroxyl stretch samples essentially every
environment within one vibrational lifetime. A static inhomogeneous distribution separable
by an echo does not exist at the relevant delay. The decomposition presupposes a
separation of timescales that this system does not supply.

Neither case is closed by simply supplying the missing number: what fails is the static two-component decomposition of Step~1 on the lifetime timescale, and a richer time-dependent treatment, not a further measurement of the same kind, is what the system requires.
Both are systems for which the scalar reduction the
procedure asks for does not have a well-formed answer, and reporting that is the
correct output. The refusal is a statement about the reduction, not a claim that
the underlying dynamics are ill defined. The distinction
between a missing measurement and an ill-posed decomposition is not visible in a
spectrum, and a procedure that returned a number in either case would be wrong
in a way no consistency check would catch.

One further observation from the same source sharpens the point, and it is
worth stating carefully because the obvious reading of it is wrong. For
W(CO)$_6$ in CHCl$_3$ the vibrational lifetime lengthens by about nine percent
between the melting and boiling points, from 322 to 350~ps, while the
absorption linewidth narrows by about nine percent over the same range, from
20.5 to 18.6~cm$^{-1}$ \citep{tokmakoff1994}. A lifetime and a width moving in
opposite senses looks at first like a contradiction, but it is not: since
$\widetilde{\Gamma}_{\mathrm{life}} = 1/2\pi cT_1$, a lengthening lifetime
predicts a narrowing lifetime width, and the two trends are therefore
directionally CONSISTENT ($-8.0$ percent predicted against $-9.3$ percent
observed). No direction test is available here, and none is claimed. What
separates the two quantities is MAGNITUDE. The lifetime-limited widths implied
by those lifetimes are 0.0165 and 0.0152~cm$^{-1}$, so the measured absorption
line is about $1.2\times10^{3}$ times broader at both ends of the range and is
dominated throughout by contributions that are not lifetime. A near-agreement
in trend across two observables separated by three orders of magnitude in size
is exactly the coincidence that makes a width look interpretable as a damping
coefficient when it is not, which is why the procedure tests the magnitude and
not the sign.
